\newpage
\section{Speedup Analysis: Baseline Deficit Attribution and Design-Parameter Sensitivity}
\label{app:deficit}

A large end-to-end speedup number is only credible if two things are shown
explicitly: (i) the naive baseline actually \emph{carries} that deficit, and
the deficit can be \emph{localized} to a physical cause; and (ii) the reported
gain can be \emph{decomposed} into a small number of design decisions, each of
which is separately sized. This appendix supplies both. We first attribute the
naive-baseline per-operation runtime into compute-active versus memory-stall
(versus communication) time (\cref{app:deficit:perop}); we then decompose the
$34.6$--$40.2\times$ co-design speedup into exactly two orthogonal factors
(\cref{app:deficit:twofactor}); we sweep each design decision one at a time and
report the objective's sensitivity (\cref{app:deficit:sweeps}); and we rank the
decisions by objective sensitivity to guide where a chip designer should spend area
and power (\cref{app:deficit:ranking}). This section is a companion to the
pre-RTL feasibility gates of \cref{app:details}: the gates certify that a design
is \emph{buildable}; this section certifies where its speedup \emph{comes from}.

\paragraph{Reading conventions (read before any number below).}
Throughout, \textbf{MFU is model-MFU}: the roofline-tightness ratio
$\mathrm{MFU}=t_{\mathrm{compute}}/t_{\mathrm{solver}}$ of \cref{app:mfu}
(cache compute floor over the LP/analytic latency within one cost-model
convention), \emph{never} a measured-silicon utilization; the searched
datapaths are hypothetical, not fabricated. \textbf{Same-silicon} comparisons
are equal-area and equal-TDP \emph{by construction} (both sides share the PE
grid, so the area/power model of \cref{app:gates} returns identical budgets).
The $97\%$ pre-fusion memory-stall figures and the naive MFU values are
\emph{naive-baseline attributions}, not properties of the co-designed schedule.
Every quantity is convention-tagged, and we never form a ratio across two
conventions. The fixed-silicon algorithmic gain on EfficientNet-B7 is
$2.21\times$ under this paper's naive convention (a $2.81\times$ value that
appears in an earlier discussion note is a transcription error and is not used).

% ---------------------------------------------------------------------------
\subsection{Per-Operation Deficit Attribution}
\label{app:deficit:perop}

To justify a custom-silicon speedup on the order of $34$--$40\times$ we must
first show that the naive baseline wastes that much time, and say \emph{where}.
We attribute the naive per-operation runtime into three buckets:
\emph{compute-active} (the systolic array is issuing useful MACs),
\emph{memory-stall} (the array is idle waiting on off-chip traffic), and, for
multi-device / mixture-of-experts execution, \emph{communication} (inter-device
transfer of dispatched tokens or expert activations). On a single package the
communication bucket is empty by construction: all eight blueprints of
\cref{app:blueprints} hold their active working set on-package (for DeepSeek-V3,
the active-expert set is resident and the routing layer, not the datapath,
pages the full model), so no operation incurs cross-device traffic.

\paragraph{The deficit is memory stall, not communication.}
Two independent attributions, on two different silicon points, agree that the
naive deficit is almost entirely memory stall.

\emph{(a) On the co-designed chip (blueprint convention).} Before fusion and
time-indexed placement, naive per-operation streaming (each operation reads
all of its tensors from DRAM and writes them back) leaves the array stalled
on memory for $\mathbf{97.1}$--$\mathbf{97.5\%}$ of naive time
(\cref{tab:deficit_perop}; individual blueprints report $97.1\%$ for
Llama-1 13B/65B, $97.2\%$ for DeepSeek-V3, $97.5\%$ for Llama-3 70B). This is
large precisely because the co-designed chip provisions $20$--$27\times$ more
MACs than TPU-v3: the compute floor shrinks, so any unfused, re-streaming
schedule becomes memory-bound. Fusion plus time-indexed residency keep the
reused tensors on-chip and recover the schedule to the compute floor,
$\mathrm{MFU}=1.000$ (0.998 on Llama-1 13B; \cref{tab:cheetah_v1_llama1_13b}).

\emph{(b) On fixed TPU-v3 silicon (grid convention).} On the $65{,}536$-MAC
baseline itself, naive streaming reaches only $\mathrm{MFU}=0.40$--$0.64$ of the
compute floor on TPU-v3 (Llama-3 70B $0.399$, DeepSeek-V3 $0.451$,
Llama-1 13B $0.467$, SD1.5-UNet $0.643$), and up to $0.78$ across the
compute-bound cells of the wider fixed-hardware grid; on the bandwidth-starved
NPU ($64$\,GB/s) it collapses to $\mathrm{MFU}=0.02$--$0.05$. The joint schedule
recovers every compute-bound cell to $\mathrm{MFU}\approx1.000$, so the
fixed-silicon speedup is a \emph{pure utilization recovery}, not additional
compute.

\begin{table}[ht!]
\centering
\caption{Naive-baseline deficit attribution. Compute-active and memory-stall
fractions are $\mathrm{MFU}_{\mathrm{naive}}$ and $1-\mathrm{MFU}_{\mathrm{naive}}$
on fixed TPU-v3 silicon (grid convention, scale-$T_{\max}$ naive). Pre-fusion
memory stall is the blueprint-chip convention (naive per-op streaming on the
co-designed datapath). Communication is zero on a single package. The last
column is the joint schedule's recovered model-MFU. All values are
naive-baseline attributions under one cost-model convention (\cref{app:mfu});
they are not measured-silicon utilizations.}
\label{tab:deficit_perop}
\small
\begin{tabular}{@{}lccccc@{}}
\toprule
& \multicolumn{3}{c}{\textbf{Naive baseline (TPU-v3, per-op)}}
& \textbf{Pre-fusion stall} & \textbf{Joint} \\
\cmidrule(lr){2-4}
\textbf{Workload} & Compute-active & Memory-stall & Comm.\ & (chip, no fusion)
& $\mathrm{MFU}$ \\
\midrule
Llama-1 13B  & 0.467 & 0.533 & 0 & 97.1\% & 1.000 \\
Llama-1 65B  & --\,$^{\dagger}$ & --\,$^{\dagger}$ & 0 & 97.1\% & 1.000 \\
Llama-3 70B  & 0.399 & 0.601 & 0 & 97.5\% & 1.000 \\
DeepSeek-V3  & 0.451 & 0.549 & 0 & 97.2\% & 1.000 \\
SD1.5-UNet   & 0.643 & 0.357 & 0 & --      & 0.963 \\
\bottomrule
\end{tabular}

\vspace{2pt}
{\footnotesize $^{\dagger}$Llama-1 65B is not one of the five workloads in the
fixed-hardware MFU grid; its pre-fusion stall is read from
\cref{tab:cheetah_v1_llama1_65b}. Across the compute-bound cells of the wider
grid, naive MFU spans $0.40$--$0.78$; the bandwidth-starved NPU cells sit at
$0.02$--$0.05$.}
\end{table}

\paragraph{The stalled time is avoidable, not compulsory.}
The $97\%$ stall is \emph{re-streaming} of tensors that could have stayed
on-chip, not irreducible DRAM traffic. Two floor analyses bound the compulsory
component. First, on each blueprint the memory floor (compulsory bytes over the
chip's own bandwidth) is only $0.10$--$0.33\%$ of the compute floor
(Llama-1 13B $2.27\times10^{7}/2.31\times10^{10}=0.10\%$; Llama-3 70B $0.31\%$;
DeepSeek-V3 $0.33\%$; \cref{app:blueprints}). Second, a mapping-independent
attainable-traffic analysis in the style of Orojenesis~\citep{huang2024mindthegap}
places the attainable per-Einsum DRAM floor at $\mathbf{0.1}$--$\mathbf{0.9\%}$
of the compute floor on every cell we measured (Llama-3 70B $0.116\%$,
DeepSeek-V3 $0.182\%$, SD1.5-UNet $0.896\%$), because at the deployed
prefill-like batch the weight reuse is large. Both bound the same conclusion:
the compulsory off-chip traffic that \emph{no} schedule can remove is under
$1\%$ of the compute floor, so essentially the entire naive memory-stall
deficit is recoverable by scheduling (fusion, time-indexed residency, and
tile choice) and the recoverable ceiling is not gated by DRAM bandwidth or
by network communication.

\paragraph{Planned figure: per-operation stacked time breakdown.}
We will add a per-operation stacked-bar figure (one bar per HLO operation,
ordered by graph position; each bar split into compute-active / memory-stall /
communication seconds) for the naive baseline and for the CHEETAH schedule side
by side, on all four LLM blueprints, in the style of the FAST per-layer
utilization plot~\citep{zhang_fast} and Eyeriss-style per-layer accounting
\citep{chen2019eyerissv2flexibleaccelerator}. The naive bars are dominated by
the memory-stall segment ($\sim97\%$ aggregate); the CHEETAH bars are almost
pure compute-active (aggregate $\mathrm{MFU}\to1$); the communication segment is
empty on a single package and appears only in the multi-device / MoE variant of
\cref{app:deficit:sweeps}. This figure is a \emph{planned appendix experiment}
(the underlying per-op compute/stall decomposition exists in the solver output;
the figure is not yet rendered).
% TODO(figure): render figs/deficit_perop_stack.pdf from the per-op
% compute/stall/comm segments; naive vs CHEETAH, four LLM blueprints.

% ---------------------------------------------------------------------------
\subsection{Two-Factor Speedup Decomposition}
\label{app:deficit:twofactor}

The headline speedup factors \emph{exactly} into two orthogonal terms:
\begin{equation}
\underbrace{\frac{T_{\mathrm{naive}}(\text{TPU-v3})}{T_{\mathrm{joint}}(\text{chip})}}_{\text{headline}}
\;=\;
\underbrace{\frac{T_{\mathrm{naive}}(\text{TPU-v3})}{T_{\mathrm{joint}}(\text{TPU-v3})}}_{f_1:\ \text{scheduling recovery}}
\;\times\;
\underbrace{\frac{T_{\mathrm{joint}}(\text{TPU-v3})}{T_{\mathrm{joint}}(\text{chip})}}_{f_2:\ \text{MAC provisioning}} .
\label{eq:twofactor}
\end{equation}
The first factor $f_1$ is a \emph{fixed-silicon}, same-area / same-TDP
algorithmic gain: it re-schedules the same $65{,}536$-MAC TPU-v3 and buys
$\mathbf{1.48}$--$\mathbf{1.78\times}$. The second factor $f_2$ holds the joint
schedule fixed and provisions the PE grid up to the co-designed chip's MAC
count, buying $\mathbf{20.0}$--$\mathbf{26.9\times}$. Their product reproduces
the per-workload Figure-2 headline (and the blueprint ``achieved speedup vs.\
naive TPU-v3'' rows) to within $1.4\%$
(\cref{tab:deficit_twofactor}); the residual is the analytic
chip-endpoint solve's utilization roll-off, cross-checked against the
blueprints' Timeloop-in-loop objective.

\begin{table}[ht!]
\centering
\caption{Two-factor decomposition of the co-design speedup
(\cref{eq:twofactor}). $f_1$ is the fixed-silicon scheduling gain on TPU-v3
(same area/TDP), $f_2$ the MAC-provisioning gain to the co-designed chip.
Reconstructed $=f_1\times f_2$; headline is the blueprint value. Convention:
config-0 naive; bf16 cycle basis; HiGHS-Optimal endpoints.}
\label{tab:deficit_twofactor}
\small
\begin{tabular}{@{}lccccc@{}}
\toprule
\textbf{Workload} & $f_1$ (sched., TPU-v3) & $f_2$ (MAC prov.)
& Reconstructed & Headline & Recon./Headline \\
\midrule
Llama-1 13B  & 1.662 & 20.31 & 33.75 & 34.22 & 0.986 \\
Llama-1 65B  & 1.727 & 20.04 & 34.60 & 34.60 & 1.000$^{\ddagger}$ \\
Llama-3 70B  & 1.484 & 26.85 & 39.84 & 40.16 & 0.992 \\
DeepSeek-V3  & 1.784 & 19.97 & 35.63 & 35.68 & 0.999 \\
\midrule
\textbf{Range} & \textbf{1.48--1.78} & \textbf{20.0--26.9}
& \multicolumn{3}{c}{\textbf{34.6--40.2} (matches Fig.~2)} \\
\bottomrule
\end{tabular}

\vspace{2pt}
{\footnotesize $^{\ddagger}$Llama-1 65B's analytic chip-endpoint solve was
degenerate under HiGHS time limits, so its $f_2$ uses the authoritative
blueprint objective (exact by construction); the other three use the
\emph{independent} analytic chip solve and land within $1.4\%$.}
\end{table}

\paragraph{Use only the two-factor form; do not multiply the full decision ladder.}
It is tempting to read the software-ablation ladder (tile, placement, remat,
datapath, datatype) as a chain of multiplicative factors and multiply them all.
That \emph{overshoots} the headline by $1.8$--$2.1\times$ (a mechanical product
of $70.7$--$73.2\times$ versus the $34$--$40\times$ headline), for two
convention reasons, neither of them physics:
\begin{enumerate}
\item \textbf{The INT8 throughput label double-counts.} The blueprint headline
$T_{\mathrm{joint}}$ is a \emph{bf16 cycle count} at the chip's MAC count; INT8
(W8A8) is carried as a throughput / power label (roughly half the area and power
per MAC), \emph{not} as a separate cycle factor. Multiplying an extra $2\times$
for INT8 into a cycle product is a double-count. Dropping it lands the mechanical
ladder at $35$--$37\times$, within $3$--$12\%$ of the headline.
\item \textbf{The software rows use a $\sim1.3\times$ smaller naive baseline}
(the scale-$T_{\max}$ naive) than the config-0 naive against which the headline
is measured, so the two families of factors are not on the same numerator.
\end{enumerate}
\Cref{eq:twofactor} is the form that reconciles \emph{exactly} because $f_1$ and
$f_2$ share a single naive numerator and a single cycle basis. As a naming
caution: the same fixed-silicon scheduling recovery is $1.48$--$1.78\times$
against the config-0 naive (used above) and $2.14$--$2.51\times$ against the
grid's scale-$T_{\max}$ naive (the fixed-hardware grid convention of
\cref{app:deficit:perop}); the two differ only by the naive-baseline
definition, and must never be crossed within one ratio.

\paragraph{The provisioning factor is sub-linear (the anti-``just add MACs''
result).} $f_2$ is not free scaling. Measured as delivered throughput per
provisioned MAC in an equal-area / equal-TDP resource-matched ablation,
$20.0$--$27.0\times$ more MACs realize only $\approx16\times$ provisioning
throughput, a linearity of $\mathbf{0.59}$--$\mathbf{0.81}$ (Llama-3 70B is
worst, $27\times$ MACs $\to$ $16.0\times$). Equivalently, the fixed-family MAC
leg from $65{,}536$ to $1{,}048{,}576$ MACs (a $16\times$ increase) yields only
$14.0$--$15.7\times$ cycle reduction, because a residual memory floor remains.
Added silicon buys sub-linearly, which is exactly why capturing the free
scheduling factor $f_1$ first (and provisioning MACs only up to the point of
diminishing return) is the efficient order of operations.

% ---------------------------------------------------------------------------
\subsection{One-Decision-at-a-Time Sensitivity Sweeps}
\label{app:deficit:sweeps}

We now hold the design fixed and sweep a single decision, reporting the objective
(LP cycles) sensitivity. Two decisions have dedicated measured sweeps (global
buffer, MAC count); the remaining fine-grained sweeps (per-unit peak-FLOPS,
interconnect radix / link count) are framed here and marked as
\textcolor{orange!85!black}{\textbf{planned}} appendix experiments so they can be
added to the paper's experiment list; we do not fabricate their numbers.

\paragraph{(a) Global-buffer size $C_{\mathrm{GM}}$: measured, near-inert in
the compute-bound regime.} \Cref{tab:deficit_cgm} sweeps $C_{\mathrm{GM}}$ from
$1$ to $64$\,MiB for DeepSeek-V3 on TPU-v3 ($900$\,GB/s), for the joint (free-$y$)
LP and for a fixed-tile FAST-subset LP. The joint LP is \emph{flat}: its objective
moves $0.39\%$ across a $64\times$ buffer sweep ($1.0723\times10^{12}\to
1.0682\times10^{12}$ cycles), and is bit-identical from $8$\,MiB upward. By
contrast the fixed-tile schedule \emph{needs} the buffer, improving $1.18\times$
over the same sweep. The joint formulation substitutes tile choice for buffer
capacity: when tiling is free, a larger global buffer buys almost nothing. The
tile-dependent capacity extension is likewise inert here: it matches stock V2
to within $0.077\%$ across the full $C_{\mathrm{GM}}\in\{4,\dots,64\}$\,MiB sweep.
Under bandwidth pressure ($64$--$128$\,GB/s) the buffer regains a modest decision
($\approx1.19\times$ for Llama-1 13B over $8\to64$\,MiB), but \emph{only through
placement}: with placement frozen (row A) the objective is constant across
$C_{\mathrm{GM}}$, so $C_{\mathrm{GM}}$ helps exclusively by enabling
time-indexed residency, never on its own. The blueprints adopt
$C_{\mathrm{GM}}=64$\,MiB.

\begin{table}[ht!]
\centering
\caption{Global-buffer sweep, DeepSeek-V3 on TPU-v3 ($900$\,GB/s, native cache
convention, HiGHS-Optimal). The joint (free-$y$) LP is essentially
$C_{\mathrm{GM}}$-insensitive ($0.39\%$ over a $64\times$ sweep); the fixed-tile
FAST-subset LP is not ($1.18\times$). Objective in LP cycles.}
\label{tab:deficit_cgm}
\small
\begin{tabular}{@{}lcccccccc@{}}
\toprule
$C_{\mathrm{GM}}$ (MiB) & 1 & 2 & 4 & 8 & 16 & 32 & 64 & sweep ($1\!\to\!64$) \\
\midrule
Joint (stock V2, free-$y$) & 1.0723 & 1.0699 & 1.0690 & 1.0682 & 1.0682
& 1.0682 & 1.0682 & $\times10^{12}$; $1.004\times$ \\
Fixed-tile (FAST-subset)   & 1.4627 & 1.4459 & 1.4327 & 1.4161 & 1.3863
& 1.3364 & 1.2371 & $\times10^{12}$; $1.182\times$ \\
\bottomrule
\end{tabular}
\end{table}

\paragraph{(b) Peak throughput / MAC-unit count: measured (count), planned
(per-unit rate).} The total-MAC provisioning sweep is measured: raising the PE
grid from $65{,}536$ to $1{,}048{,}576$ MACs (a $16\times$ increase) reduces the
objective by $14.0$--$15.7\times$ (Llama-1 13B $15.72\times$, DeepSeek-V3
$14.98\times$, Llama-3 70B $14.04\times$, Llama-1 65B $14.02\times$), and the
provisioning factor to the full co-designed chip is $f_2=20.0$--$26.9\times$
(\cref{app:deficit:twofactor}), both sub-linear (\S\ref{app:deficit:twofactor}).
This is the single largest decision on the objective. A \emph{finer} sensitivity
that varies per-unit peak throughput (systolic-array dimensions, vector-lane
width, or clock) at a \emph{fixed} MAC count and fixed area, separating
``more units'' from ``faster units'', is not yet run and is marked
\textcolor{orange!85!black}{\textbf{planned}}; the datapath template of
\cref{app:gates} and MAGNet-style generators~\citep{venkatesan2019magnet}
already parameterize it.

\paragraph{(c) DRAM bandwidth: measured, inert in the compute-bound regime.}
Sweeping bandwidth from $900$ to $3{,}350$\,GB/s at the fixed $65{,}536$-MAC
point moves the objective by $\le0.25\%$ (BW leg $=1.000$--$1.0025\times$ for
all four LLMs): at the deployed prefill-like batch these workloads are
compute-bound, so bandwidth beyond the floor buys nothing. Bandwidth becomes the
dominant decision only in the bandwidth-starved regime: at $64$\,GB/s (the compact
NPU) the joint-vs-naive fixed-hardware gain is $5.9$--$31.4\times$, because there
the roofline is memory-bound and residency / placement recover a much larger
share. Thus bandwidth sensitivity is regime-dependent: near-zero when
compute-bound, first-order when memory-bound.

\paragraph{(d) Interconnect: on-chip NoC and inter-device topology (planned).}
Two interconnect decisions are framed but not yet swept.
\emph{On-chip}, the NoC / crossbar radix and link count set the PE-to-SRAM
routing capacity; this is the resource bounded by the crossbar / NoC-routing
gate of \cref{app:gates} (gate 4, $N_{\mathrm{NoC}}=n_{\mathrm{PEs}}\times
n_{\mathrm{banks}}\le5\times10^{6}$ cells), and its provisioning follows the
network-on-package accounting of Simba~\citep{shao2019simba}. In the single-chip
compute-bound regime the on-chip NoC is provisioned to the compute floor and is
not the binding resource, but a dedicated radix / link sweep (objective
versus NoC cells at fixed MACs) is
\textcolor{orange!85!black}{\textbf{planned}}.
\emph{Across devices}, in the multi-GPU / mixture-of-experts regime the
inter-device topology (radix and link count of the NVLink-class fabric) binds
only when experts are dispatched \emph{across} devices; when the active-expert
set is package-resident (as in every blueprint here) the communication term is
zero (\cref{app:deficit:perop}). A distributed / MoE topology sweep (objective
and the communication bucket versus inter-GPU radix and link count, in the
spirit of Chiplet-Cloud and wafer-scale serving
analyses~\citep{peng2024chipletcloudbuildingai,zhu2024theseusexploringefficientwaferscale}
and co-optimized network provisioning~\citep{rashidi2023unico}) is
\textcolor{orange!85!black}{\textbf{planned}}.

\begin{table}[ht!]
\centering
\caption{Single-decision sensitivity summary (compute-bound LLM regime, TPU-v3
operating point). ``Objective sensitivity'' is the LP-cycle response to the
stated sweep. Cells marked \textcolor{orange!85!black}{\textbf{planned}} have no
dedicated run yet and are listed as appendix experiments; their numbers are not
fabricated.}
\label{tab:deficit_sweepsummary}
\small
\begin{tabular}{@{}llll@{}}
\toprule
\textbf{Decision} & \textbf{Sweep} & \textbf{Objective sensitivity} & \textbf{Status} \\
\midrule
Global buffer $C_{\mathrm{GM}}$ & $1\to64$\,MiB (BW$=900$)
& $1.004\times$ (joint); $1.18\times$ (fixed-tile) & measured \\
& under BW pressure ($\le128$) & $\approx1.19\times$, via placement only & measured \\
MAC count (total) & $65{,}536\to1{,}048{,}576$ ($16\times$)
& $14.0$--$15.7\times$ (sub-linear); $f_2=20.0$--$26.9\times$ & measured \\
Peak rate / unit & systolic dims, clock @ fixed MACs
& -- & \textcolor{orange!85!black}{\textbf{planned}} \\
DRAM bandwidth & $900\to3{,}350$\,GB/s & $\le1.003\times$ (compute-bound) & measured \\
& $64$\,GB/s (NPU) & $5.9$--$31.4\times$ (memory-bound) & measured \\
On-chip NoC radix / links & $N_{\mathrm{NoC}}$ at fixed MACs
& -- (gate 4 bounds it) & \textcolor{orange!85!black}{\textbf{planned}} \\
Inter-device topology (MoE) & inter-GPU radix / links
& -- (binds on cross-device dispatch) & \textcolor{orange!85!black}{\textbf{planned}} \\
\bottomrule
\end{tabular}
\end{table}

% ---------------------------------------------------------------------------
\subsection{Decision-Impact Ranking}
\label{app:deficit:ranking}

\Cref{tab:deficit_ranking} ranks the decisions by objective sensitivity in the
compute-bound LLM regime and states \emph{why} each sits where it does, so a
chip designer can read off where area and power are best spent.

\begin{table}[ht!]
\centering
\caption{Design decisions ranked by objective sensitivity in the compute-bound LLM
regime (TPU-v3 operating point). ``Cost'' is what the decision spends on the die.
The ranking inverts in the memory-bound / distributed regimes, as noted.}
\label{tab:deficit_ranking}
\small
\begin{tabular}{@{}clllp{3.5cm}@{}}
\toprule
\textbf{Rank} & \textbf{Decision} & \textbf{Sensitivity} & \textbf{Cost}
& \textbf{Binds when} \\
\midrule
1 & MAC-count provisioning & $20.0$--$26.9\times$ (sub-linear) & area $+$ power
& always (compute-bound); diminishing returns \\
2 & Scheduling (tile+fusion+placement) & $1.48$--$1.78\times$ & none (same silicon)
& always; also a feasibility enabler at tight buffers \\
3 & DRAM bandwidth & $\approx1.0\times$; $5.9$--$31.4\times$ if starved & pins $+$ power
& memory-bound regime only \\
4 & Global buffer $C_{\mathrm{GM}}$ & $\approx1.0\times$; $\le1.19\times$ pressed & area (SRAM)
& buffer-pressed regime; only via placement \\
5 & Interconnect (NoC / inter-GPU) & $\approx1.0\times$ on-chip & area $+$ links
& distributed / MoE cross-device dispatch \\
\bottomrule
\end{tabular}
\end{table}

\paragraph{Why this ordering, and what it tells a designer.}
In the compute-bound LLM regime the objective is set by how many useful MACs the
schedule can keep busy. \textbf{(1) MAC-count provisioning} moves the objective
most ($20$--$27\times$) but sub-linearly and at direct area/power cost, so it is
the decision to spend on, up to the point of diminishing return, and never past
the feasibility gates of \cref{app:gates} (a bigger matrix unit can even degrade
delivered performance, as observed for production TPU
silicon~\citep{jouppi2017datacenter}). \textbf{(2) Scheduling} is a
\emph{free} multiplier ($1.48$--$1.78\times$ at equal area and TDP) that must be
captured first: it costs no silicon, it is what recovers naive MFU $0.40$ to
$\approx1.0$, and at tight buffers it is a feasibility enabler, not merely an
optimization (on EfficientNet-B7 at TPU-v3's $32$\,MiB buffer the pure-streaming
row is infeasible and time-indexed placement restores feasibility: placement
alone at $2.03\times$, joint optimum $2.21\times$). \textbf{(3) DRAM bandwidth} and \textbf{(4) global buffer} are
near-inert once compute-bound (bandwidth beyond the floor and buffer beyond a
few MiB buy essentially nothing; \cref{tab:deficit_cgm}), and the joint LP
actively trades buffer for tile choice, so paying for HBM or a large global SRAM
buys lower TDP-per-TFLOP, not more throughput, at these operating points. Both
become first-order only when the workload is genuinely memory-bound (low
bandwidth, or decode-time batch-1 operational intensity). \textbf{(5)
Interconnect} is inert on a single compute-bound chip and matters only in the
distributed / MoE regime, when experts are dispatched across devices and the
communication bucket becomes nonzero. The practical guidance: \emph{buy MACs
(watching the sub-linear roll-off and the gates), capture scheduling for free,
and invest in bandwidth, buffer, or interconnect only when the regime (memory-
or communication-bound) actually makes them bind.}

\paragraph{Sweeps requiring new runs (for the paper's experiment list).}
The following are framed above but not yet run, and are proposed as appendix
experiments: (i) a per-unit peak-throughput sweep (systolic dimensions / clock
at fixed MAC count and area), isolating ``faster units'' from ``more units'';
(ii) an on-chip NoC radix / link-count sweep (objective versus $N_{\mathrm{NoC}}$
at fixed MACs, tied to gate 4); (iii) an inter-device / MoE topology sweep
(objective and the communication bucket versus inter-GPU radix and link count,
with experts dispatched across devices); and (iv) the per-operation stacked
compute/stall/communication time-breakdown figure of \cref{app:deficit:perop}.
